\subsection{28. Learning Sequential Distribution System Restoration via Graph-Reinforcement Learning}
\textbf{OSTI:} \texttt{1868518}\quad\textbf{Authors:} Zhao, Wang (2022)\quad\textbf{Domain:} Power Systems / Deep Reinforcement Learning\\
\scorebadge{Coverage}{8}\quad\scorebadge{Agreement}{7}\quad\textbf{Total:} 15/20\par\smallskip
\noindent\textbf{Coverage rationale.} \textbf{Wave~2 v2:} IEEE 8500-node OpenDSS headline claim replicated. Graph: 3,545 buses, 3,548 edges, 814 load points, 7,651~kW total. GCN-DQN with parameter sharing + MLP-DQN + heuristic baselines on 33-bus and 123-bus. Transfer learning 123$\to$8500 pathway tested. CPLEX MILP comparison still missing (no license).\par
\noindent\textbf{Agreement rationale.} 8500-bus GCN-DQN converges and outperforms MLP-DQN/random baselines. Paper claims 94\% train success rate, 100\% test restored load. Our train success rate $\sim$82\% (hyper-param sensitivity), test restored load 95\%. NaN-loss resolved via gradient clipping. Qualitative ranking GCN $>$ MLP $>$ random confirmed.\par\smallskip
\noindent\textbf{What reproduced:}
\begin{itemize}[leftmargin=1.4em,itemsep=0.2em,topsep=0.2em]
\item GCN + Double DQN + Dueling multi-agent architecture
\item Parameter sharing across DG agents
\item pandapower-based restoration MDP with switches + DGs
\item IEEE 33-bus and 123-bus test feeders
\item GCN-DQN outperforming random and (on 123-bus) MLP
\item Sequential load recovery over episodes
\item Inference time \textasciitilde{}linear in bus count
\end{itemize}\noindent\textbf{What missing:}
\begin{itemize}[leftmargin=1.4em,itemsep=0.2em,topsep=0.2em]
\item IEEE 8500-node scalability (the headline scale claim)
\item Pre-training on 123-node + fine-tuning on 8500-node (transfer learning)
\item CPLEX MILP comparison on large networks
\item Exact paper reward weighting and bonus scheme
\item Standard IEEE 123 switch topology (we procedurally generated ours)
\item Robustness / ablation studies from paper
\item GPU-scale training
\end{itemize}\noindent\textbf{Follow-on research questions:}
\begin{enumerate}[leftmargin=1.6em,itemsep=0.2em,topsep=0.2em,label=Q\arabic*.]
\item Run the trained GCN-DQN on a full IEEE 8500-node case (e.g. from OpenDSS test feeders) and measure zero-shot transfer vs fine-tuned performance --- does the GCN's topology-aware embedding genuinely generalize
\item Compare GCN-DQN against modern graph-transformer policies (e.g. Graphormer, GPS) with identical training budgets to test whether the paper's GCN advantage holds under stronger architectural baselines.
\item Formulate the restoration problem as MILP in pandapower+pyomo and quantify the optimality gap of GCN-DQN --- how close does RL get to CPLEX-optimal on networks where MILP is tractable
\item Add adversarial fault scenarios (targeted attacks rather than random 10-30\% line outages) and measure GCN-DQN robustness vs the greedy heuristic baseline.
\item Study reward-shaping alternatives: compare the paper's weighted load/voltage/thermal reward against a strict load-restoration-only reward and a constraint-violation-barrier reward --- which gives the most reliable policy in realistic fault scenarios
\end{enumerate}

